% Este trabalho está licenciado sob a licença Creative Commons CC-BY-SA.
% Você é livre para usar, compartilhar e adaptar este material, desde que seja atribuída a autoria e que as obras derivadas mantenham a mesma licença.

% ---
% Pacotes Principais
% ---
\usepackage[left=3cm,right=2cm,top=3cm,bottom=2cm]{geometry}    % layout de página
\usepackage{lmodern}		% Usa a fonte Latin Modern
\usepackage[T1]{fontenc}	% Selecao de codigos de fonte.
\usepackage[utf8]{inputenc}	% Codificacao do documento (conversão automática dos acentos)
\usepackage{indentfirst}	% Indenta o primeiro parágrafo de cada seção.
\usepackage{color}			% Controle das cores
\usepackage{xcolor}
\usepackage{graphicx}		% Inclusão de gráficos
\usepackage{microtype} 		% para melhorias de justificação
\usepackage{pslatex}        % torna o documento padrão para as fontes Times/Helvetica/Courier
% \usepackage{url}
\usepackage{booktabs}
\usepackage{subfig}
\usepackage{pdfpages}
\usepackage{lastpage}
\usepackage{fancyhdr}
\usepackage{float}
\usepackage{caption}
\usepackage{amsmath, amssymb, amsthm, amsfonts}
\usepackage{setspace}
\usepackage{chngcntr}
\usepackage{keystroke}
\usepackage{tocloft}
\usepackage{longtable, ltcaption}
\usepackage{newfloat}
\usepackage{tikz}
\usepackage{hyperref}
\usepackage{fancyvrb}
\usepackage{titlesec}
\usepackage{ifthen}
\usepackage{tabularx}
\usepackage{algorithm}
\usepackage{algpseudocode}
\usepackage{amsmath}
\usepackage{pgfgantt}
% ---
% ---
% Pacote de citações
% ---
% \usepackage[alf,abnt-etal-text=it,abnt-repeated-author-omit=yes,abnt-etal-list=0,abnt-etal-cite=4,abnt-emphasize=bf]{abntex2cite}
\usepackage[alf,abnt-etal-text=it,abnt-repeated-author-omit=no,abnt-etal-list=0,abnt-etal-cite=4,abnt-emphasize=bf]{abntex2cite}
% ---

\usepackage{listings}
\lstset{
  xleftmargin=1cm,      % 1 cm de "indentação" à esquerda
  breaklines=true,      % quebra linhas longas
  breakatwhitespace=true,
  basicstyle=\ttfamily\small
}

% Troca "Listing" por "Listagem"
\renewcommand{\lstlistingname}{Listagem}

% Se quiser também mudar o plural:
\renewcommand{\lstlistlistingname}{Lista de Listagens}

% Define enumerate no padrão a), b), c)...
\setlist[enumerate]{label=\alph*), leftmargin=3em}
\setlist[itemize]{leftmargin=3em}

\newcommand{\sourceFig}[1]{%
    \vspace{-2.5ex}
    \caption*{\small Fonte: #1}%
}

\newcommand{\sourceTab}[1]{%
    \caption*{\small Fonte: #1}%
}



% ---
% Customizações de Fontes e Tamanhos de Letras usados na Classe ABNTEX2
% ---
\renewcommand{\chaptitlefont}{\ABNTEXchapterfont\ABNTEXchapterfontsize}

\setsecheadstyle{\ABNTEXsectionfont\ABNTEXsectionfontsize\ABNTEXsectionupperifneeded}

\renewcommand{\ABNTEXchapterfont}{\rmfamily}
\renewcommand{\ABNTEXchapterfontsize}{\normalsize\bfseries}

\renewcommand{\ABNTEXpartfont}{\ABNTEXchapterfont}
\renewcommand{\ABNTEXpartfontsize}{\LARGE\bfseries}
\renewcommand{\cftpartfont}{\normalfont\rmfamily\bfseries}
\renewcommand{\cftpartpagefont}{\normalfont\rmfamily\bfseries}

\renewcommand{\ABNTEXsectionfont}{\ABNTEXchapterfont}
\renewcommand{\ABNTEXsectionfontsize}{\normalsize}

\titleformat{\section}
  {\MakeUppercase}   % transforma em caixa alta
  {\thesection}{1em}{}
  
\renewcommand{\ABNTEXsubsectionfont}{\ABNTEXsectionfont}
\renewcommand{\ABNTEXsubsectionfontsize}{\normalsize}


\titleformat{\subsection}
  {\bfseries}   % transforma em negrito
  {\thesubsection}{1em}{}


\renewcommand{\ABNTEXsubsubsectionfont}{\ABNTEXsubsectionfont}
\renewcommand{\ABNTEXsubsubsectionfontsize}{\normalsize}

\renewcommand{\ABNTEXsubsubsubsectionfont}{\ABNTEXsubsectionfont}
\renewcommand{\ABNTEXsubsubsubsectionfontsize}{\normalsize}
% ---

% ---
% Customização do Sumário (Distância entre número e nome do item)
% ---
\setlength{\cftlastnumwidth}{3.5em}
\cftsetindents{part}{0em}{\cftlastnumwidth}
\cftsetindents{chapter}{0em}{\cftlastnumwidth}
\cftsetindents{section}{0em}{\cftlastnumwidth}
\cftsetindents{subsection}{0em}{\cftlastnumwidth}
\cftsetindents{subsubsection}{0em}{\cftlastnumwidth}
\cftsetindents{paragraph}{0em}{\cftlastnumwidth}
\cftsetindents{subparagraph}{0em}{\cftlastnumwidth}


% entries in the table of contents
\renewcommand{\cftsectionfont}{\normalfont\MakeUppercase}
\renewcommand{\cftsectionpagefont}{\normalfont}

% subsection entries (in the table of contents)
\renewcommand{\cftsubsectionfont}{\bfseries}
\renewcommand{\cftsubsectionpagefont}{\bfseries}

% subsection entries (in the table of contents)
\renewcommand{\cftsubsubsectionfont}{\normalfont}
\renewcommand{\cftsubsubsectionpagefont}{\normalfont}

\captionsetup{justification=centering, font=small}

   
% ---
% Possibilita criação de Quadros e Lista de quadros.
% Ver https://github.com/abntex/abntex2/issues/176
%
\renewcommand{\insertchapterspace}{}
\newcommand{\listofquadrosname}{Lista de Quadros}
\DeclareFloatingEnvironment[
    fileext=loq,
    listname={Lista de Quadros},
    name=Quadro,
    placement=p,
    within=none,
    chapterlistsgaps=off,
    ]{quadro}

\newlistof{listofquadros}{loq}{\listofquadrosname}
\newlistentry{quadro}{loq}{0}

% configurações para atender às regras da ABNT
\setfloatadjustment{quadro}{\centering}
\counterwithout{quadro}{chapter}
\renewcommand{\cftquadroname}{\quadroname\space} 
\renewcommand*{\cftquadroaftersnum}{\hfill--\hfill}


\setfloatlocations{quadro}{hbtp} % Ver https://github.com/abntex/abntex2/issues/176
% ---



% ---
% Possibilita criação de Gráficos e Lista de Gráficos.
% Ver https://github.com/abntex/abntex2/issues/176
%
\renewcommand{\insertchapterspace}{}
\newcommand{\listofplotsname}{Lista de Gráficos}
\DeclareFloatingEnvironment[
    fileext=lop,
    listname={Lista de Gráficos},
    name=Gráfico,
    placement=p,
    within=none,
    chapterlistsgaps=off,
    ]{plot}

\newlistof{listofplots}{lop}{\listofplotsname}
\newlistentry{plot}{lop}{0}

% configurações para atender às regras da ABNT
\setfloatadjustment{plot}{\centering}
\counterwithout{plot}{chapter}
\renewcommand{\cftplotname}{\plotname\space} 
\renewcommand*{\cftplotaftersnum}{\hfill--\hfill}


\setfloatlocations{plot}{hbtp} % Ver https://github.com/abntex/abntex2/issues/176
% ---





\setlength{\cftbeforechapterskip}{0pt} % Altera o espaçamento entre itens do sumário
% ---

% ---
% Customizações Adicionais
% ---


\captionsetup[figure]{aboveskip=2pt, belowskip=0pt}
\captionsetup[plot]{aboveskip=2pt, belowskip=0pt}
\captionsetup[table]{aboveskip=2pt, belowskip=0pt}
\captionsetup[quadro]{aboveskip=2pt, belowskip=0pt}







\providecommand{\imprimirEscola}{}
\newcommand{\escola}[1]{\renewcommand{\imprimirEscola}{#1}}

\providecommand{\imprimirCurso}{}
\newcommand{\curso}[1]{\renewcommand{\imprimirCurso}{#1}}

\providecommand{\imprimirCidade}{}
\newcommand{\cidade}[1]{\renewcommand{\imprimirCidade}{#1}}

\providecommand{\imprimirsubtitulo}{}
\newcommand{\subtitulo}[1]{\renewcommand{\imprimirsubtitulo}{#1}}

\providecommand{\imprimirexaminadorA}{}
\newcommand{\examinadorA}[1]{\renewcommand{\imprimirexaminadorA}{#1}}

\providecommand{\imprimirexaminadorB}{}
\newcommand{\examinadorB}[1]{\renewcommand{\imprimirexaminadorB}{#1}}

\counterwithout{equation}{chapter} % define uma numeração contínua (1,2,3, ...) de equações. Para uma numeração com dependência do número do capítulo (1.1, 1.2, ...), basta apagar esta linha.

\setlength{\ABNTEXsignwidth}{10cm} % define o comprimento da linha de assinatura na folha de aprovação

\setlength{\ABNTEXsignthickness}{0.5pt} % define a espessura da linha de assinatura na folha de aprovação

% Redefine o ambiente citacao para que as citações diretas com mais de 3 linhas tenham recuo de 4 cm
\renewenvironment{citacao}[1][default]{%
   \list{}{\leftmargin=0cm}%
   %\ABNTEXfontereduzida%
   \small % 10pt
   \addtolength{\leftskip}{\ABNTEXcitacaorecuo} % recuo de 4 cm da margem esquerda
   \item[]%
   \begin{SingleSpace}%
   \ifthenelse{\not\equal{#1}{default}}{\itshape\selectlanguage{#1}}{}%
 }{%
   \end{SingleSpace}%
   \endlist}%


% O tamanho do parágrafo é dado por:
\setlength{\parindent}{1.25cm}

%\setlength{\parskip}{0.2cm}  % tente também \onelineskip

% Controle do espaçamento entre um parágrafo e outro:
\setlength{\cftparskip}{0pt}

%\renewcommand{\baselinestretch}{1.5}
\linespread{1.5} % espaço entre linhas 1,5cm

% Altera o aspecto da cor azul
\definecolor{blue}{RGB}{41,5,195}
% ---


\newtheorem{definition}{Definição}[chapter]
\newtheorem{theorem}{Teorema}[chapter]
\newtheorem{corollary}{Corolário}[theorem]
\newtheorem{lemma}[theorem]{Lema}

% informações do PDF
\makeatletter
\hypersetup{
     	%pagebackref=true,
		pdftitle={\@title}, 
		pdfauthor={\@author},
    	pdfsubject={\imprimirpreambulo},
	    pdfcreator={LaTeX with abnTeX2},
		pdfkeywords={abnt}{latex}{abntex}{abntex2}{projeto de pesquisa}, 
		colorlinks=true,       		% false: boxed links; true: colored links
    	linkcolor=black,          	% color of internal links
    	citecolor=black,        		% color of links to bibliography
    	filecolor=black,      		% color of file links
		urlcolor=black,
		bookmarksdepth=4
}
\makeatother
% --- 

% ---
% Modelo de Capa conforme Manual PUCPR
\renewcommand{\imprimircapa}{
\begin{capa}
    \begin{center}
        % \includegraphics[height=15ex]{img/brasao.png}\\[1.5ex]
 		\MakeUppercase{\textbf{\imprimirinstituicao}}\\
 		\MakeUppercase{\textbf{\imprimirEscola}}\\
 		\MakeUppercase{\textbf{\imprimirCurso}}\\
 		\vspace{1.3cm}
		\MakeUppercase{\textbf{\imprimirautor}}\\
        \vspace{7cm}
        \MakeUppercase{\textbf{\imprimirtitulo}}\\
        \MakeUppercase{\textbf{\imprimirsubtitulo}}
        \vfill
        \MakeUppercase{\textbf{\imprimirCidade}}\\
        \textbf{\imprimirdata}
    \end{center}
\end{capa}
}

% ---
% Modelo de Folha de Rosto conforme Manual PUCPR
\renewcommand{\imprimirfolhaderosto}{

% =========================
% Preâmbulo automático
% Idealmente não deve ser alterado manualmente.
% Apenas as variáveis em dados_academicos.tex devem ser alteradas.
% =========================

\newcommand{\textopreambulo}{}

\ifthenelse{\equal{\imprimirTrabalhoTipo}{tcc}}{
  \renewcommand{\textopreambulo}{
    Trabalho de Conclusão de Curso apresentado ao Curso de Graduação em \imprimirCurso\ da \imprimirinstituicao,
    como requisito parcial à obtenção do título de \imprimirTituloConcedido.
  }
}{}

\ifthenelse{\equal{\imprimirTrabalhoTipo}{espec}}{
  \renewcommand{\textopreambulo}{
    Trabalho de Conclusão de Curso apresentado ao Curso de Especialização em \imprimirCurso\ da \imprimirinstituicao,
    como requisito parcial à obtenção do título de \imprimirTituloConcedido.
  }
}{}

\ifthenelse{\equal{\imprimirTrabalhoTipo}{msc-qual}}{
  \renewcommand{\textopreambulo}{
    Documento apresentado ao \imprimirPrograma\ da \imprimirinstituicao,
    como requisito parcial à obtenção do título de \imprimirTituloConcedido.
  }
}{}

\ifthenelse{\equal{\imprimirTrabalhoTipo}{msc}}{
  \renewcommand{\textopreambulo}{
    Dissertação apresentada ao \imprimirPrograma\ da \imprimirinstituicao,
    como requisito parcial à obtenção do título de \imprimirTituloConcedido.
  }
}{}

\ifthenelse{\equal{\imprimirTrabalhoTipo}{phd-qual}}{
  \renewcommand{\textopreambulo}{
    Documento apresentado ao \imprimirPrograma\ da \imprimirinstituicao,
    como requisito parcial à obtenção do título de \imprimirTituloConcedido.
  }
}{}

\ifthenelse{\equal{\imprimirTrabalhoTipo}{phd}}{
  \renewcommand{\textopreambulo}{
    Tese apresentada ao \imprimirPrograma\ da \imprimirinstituicao,
    como requisito parcial à obtenção do título de \imprimirTituloConcedido.
  }
}{}


% Linha opcional do coorientador (se não estiver vazio)
\newcommand{\linhaCoorientador}{}
\edef\tempCoorientador{\imprimirNomeCoorientador}
\if\relax\detokenize\expandafter{\tempCoorientador}\relax
  % vazio, não faz nada
\else
  \renewcommand{\linhaCoorientador}{\\ Coorientador: \imprimirNomeCoorientador}
\fi

\preambulo{
  \textopreambulo\\
  Orientador: \imprimirNomeOrientador
  \linhaCoorientador
}

\begin{folhaderosto}
    \begin{center}
     	\MakeUppercase{\textbf{\imprimirautor}}\\
     	\vspace{10.6cm}
     	\MakeUppercase{\textbf{\imprimirtitulo}} \\
     	\MakeUppercase{\textbf{\imprimirsubtitulo}}
     \end{center}
     \vspace{0.4cm}
     \hspace{8cm} % recuo de 8cm
        \begin{minipage}{8cm}
			\SingleSpacing % espaçamento simples
			\imprimirpreambulo
		\end{minipage}\\
	\begin{center}
        \vfill
        \MakeUppercase{\textbf{\imprimirCidade}}\\
        \textbf{\imprimirdata}
    \end{center}
\end{folhaderosto}
}
% ---

\makeatletter
\renewcommand{\cite}[2][]{%
  \begingroup
    \edef\tempa{#1}%
    % Sem página
    \ifx\tempa\@empty
      (\citeauthoronline{#2},~\citeyear{#2})%
    % Com página
    \else
      (\citeauthoronline{#2},~\citeyear{#2},~#1)%
    \fi
  \endgroup
}
\makeatother

\renewcommand{\citeauthor}[1]{\citeauthoronline{#1}}

% Ajustar o ambiente agradecimentos
\makeatletter
\renewenvironment{agradecimentos}{%
  \chapter*{Agradecimentos}%
  % sobe o texto em aproximadamente uma linha
  \vspace*{-\baselineskip}%
  \begingroup
    \setlength{\parindent}{0pt}% como você colocou no AtBeginEnvironment
    \setlength{\parskip}{0pt}% sem espaço antes do primeiro parágrafo
    % a partir do segundo parágrafo, usar \baselineskip como parskip
    \let\oldpar\par
    \def\par{\oldpar\setlength{\parskip}{\baselineskip}}%
}{%
  \par
  \endgroup
  \clearpage
}
\makeatother

% Menos hifenizações, mas ainda permitindo quando necessário
\sloppy
\pretolerance=2000
\tolerance=3000
\hyphenpenalty=600
\exhyphenpenalty=600

\setlist[itemize]{leftmargin=1.06em, label=\textendash}

\usepackage{caption}
\usepackage{float} % para [H]

\newenvironment{figura}[2][H]{%
  \begin{figure}[#1]
    \centering
    \begin{minipage}{#2}
      \captionsetup{width=\textwidth,justification=raggedright,singlelinecheck=false}%
}{%
    \end{minipage}
  \end{figure}
}

\newenvironment{tabela}[2][H]{%
  \begin{table}[#1]
    \centering
    \begin{minipage}{#2}
      \captionsetup{width=\textwidth,justification=raggedright,singlelinecheck=false}%
}{%
    \end{minipage}
  \end{table}
}

% Novo ambiente: quadro2
\newenvironment{quadro2}[2][H]{%
  \begin{quadro}[#1]
    \centering
    \begin{minipage}{#2}
      \captionsetup{width=\textwidth,justification=raggedright,singlelinecheck=false}%
}{%
    \end{minipage}
  \end{quadro}
}

\newenvironment{grafico}[2][H]{%
  \begin{plot}[#1]
    \centering
    \begin{minipage}{#2}
      \captionsetup{width=\textwidth,justification=raggedright,singlelinecheck=false}%
}{%
    \end{minipage}
  \end{plot}
}

\renewcommand\url[1]{#1}

% =========================
% Setters + getters (imprimir)
% =========================
\makeatletter
\newcommand{\trabalhoTipo}[1]{\def\@trabalhoTipo{#1}}
\newcommand{\tituloConcedido}[1]{\def\@tituloConcedido{#1}}
\newcommand{\programa}[1]{\def\@programa{#1}}
\newcommand{\nomeOrientador}[1]{\def\@nomeOrientador{#1}}
\newcommand{\nomeCoorientador}[1]{\def\@nomeCoorientador{#1}}

% Getters (sem @) para usar no texto com segurança
\newcommand{\imprimirTrabalhoTipo}{\@trabalhoTipo}
\newcommand{\imprimirTituloConcedido}{\@tituloConcedido}
\newcommand{\imprimirPrograma}{\@programa}
\newcommand{\imprimirNomeOrientador}{\@nomeOrientador}
\newcommand{\imprimirNomeCoorientador}{\@nomeCoorientador}
\makeatother